\SourceLeaf{385}%
\textsc{Zach} im Jahre 1809 dort ausgeführte Triangulirung betreffen, noch mehr ins Licht gesetzt werden.

Im folgenden Abschnitt findet man verschiedene mit dem Zustande der Atmosphäre im Zusammenhange stehende interessante Beobachtungen und Untersuchungen, nemlich gleichzeitige meteorologische Beobachtungen im Hospiz des Mont Cenis und in Mailand; barometrische Höhenbestimmung des ersteren Punktes und des Mont Colombier; trigonometrische Höhenbestimmung des Mont-blanc (4802,7 Meter) und des Monte Rosa (4619,6 Meter); endlich Untersuchungen über die terrestrische Refraction. Letztere wird aus den an drei Punkten (Mailand, Turin, Mondovi) beobachteten Elevationen dreier Berge Rocca Melone, Monte Viso und Monte Rosa bestimmt, wobei die absoluten Höhen der drei Standpunkte vorausgesetzt, und die Höhen der beobachteten Punkte eliminirt werden, ein Verfahren, welches wenig Sicherheit geben kann, und wie eine genauere Prüfung zeigt, auch wenig Übereinstimmung gegeben hat. Wir möchten also auf das Endresultat für das Verhältniss der Erdkrümmung zur \textit{ganzen} Refraction (5,28 zu 1) wenig Gewicht legen: die Berechnung von sechs Paaren reciproker Zenithdistanzen zwischen Hauptdreieckspunkten gibt uns dieses Verhältniss im Mittel weit kleiner, nemlich 1 zu 0,1235, sehr nahe übereinstimmend mit den bei der Hannoverschen und Liefländischen Gradmessung gefundenen Resultaten.

Der zehnte Abschnitt beschäftigt sich mit der vielbehandelten Aufgabe, aus der Breite eines Endpunktes einer gegebenen Dreiecksseite und deren Azimuth in jenem Endpunkte, dieselben Dinge für den andern Endpunkt, und den Längenunterschied auf dem elliptischen Sphäroid zu finden. Die Entwickelung enthält nur eine Umformung der \textsc{Legendre}schen Formeln, um anstatt der sogenannten reducirten Breite die wahre einzuführen. Die Formeln sind bis zu den Grössen der dritten Ordnung genau, insofern man die Abplattung und die Dreiecksseite (den Erdradius als Einheit angenommen) wie Grössen der ersten Ordnung betrachtet. Bei der Anwendung auf die gegenwärtigen Messungen hat man die Grössen der dritten Ordnung weggelassen, weil diess für die Ausübung genau genug sei. Diess ist jedoch nur insofern zuzugeben, als man die Resultate bloss zur Vergleichung mit den astonomischen Bestimmungen gebrauchen will, wo es allerdings unnöthig ist, in die Berechnung von jenen eine viel grössere Schärfe zu legen, als diese zulassen. Geht man aber von einem andern Gesichtspunkte aus, nemlich die geodaetischen Resultate so genau zu berechnen, wie es die Mes%
